\section{Legacy WebGPU Chapter Material}
\label{ch:webgpu-backend}

\texttt{WEBGPU} is CNA's fifth graphics renderer, and its youngest: the project owner
explicitly lifted a former prohibition on this renderer and authorized its implementation.
It targets native \texttt{wgpu-native} (pinned to version 29.0.1.1), selected the same way
as every other renderer via \texttt{-DCNA\_GRAPHICS\_RENDERER=WEBGPU}. Its own tracking plan
spans fourteen phases and over 130 individually numbered tasks --- large enough that this
chapter covers its trajectory rather than every task.

\section{What ``experimental'' concretely means here}

Unlike the four renderers covered in the previous three chapters, \texttt{WEBGPU} is
deliberately excluded from the cross-renderer feature matrix used throughout
Chapters~\ref{ch:sdl-renderer}--\ref{ch:vulkan-backend}, on the matrix's own stated grounds
that its feature surface was not yet broad enough for a meaningful side-by-side comparison at
the time that matrix was written. This chapter instead draws on the renderer's own dedicated
documentation and milestone history.

\section{The native 2D baseline}

Early milestones established device/surface setup, clear/present, a working
\cnaclass{Texture2D} path, vertex/index buffer uploads, and a WGSL-based
\cnaclass{SpriteBatch} implementation. A 120-frame clear-and-present lifecycle test, and a
manually-reviewed \cnaclass{SpriteBatch} validation scene (covering cropping, tint/alpha,
rotation, flips, filtering, and all three address modes), both passed before the native 2D
baseline was declared closed. The most instructive milestone here, though, is an independent
one: a real, independently-developed XNA-style game (a \emph{Windows Phone Speedy Blupi}
port, from the \texttt{mobile-eggbert} sibling repository) was run against this renderer on a
real desktop, reached its main menu, rendered its \cnaclass{SpriteBatch} content
pixel-correctly, and played through a real animated cutscene --- with progress bar and
cross-fade --- triggered by a simulated button click, with no WebGPU validation errors of any
kind. This is a genuinely different, and stronger, kind of proof than any of this renderer's
own purpose-built test scenes: an unmodified, independently-authored game simply working.

\section{A bug a manual screenshot review missed, and a pixel assertion caught}

One specific finding is worth naming as a methodological lesson in its own right: the
SpriteBatch pipeline's blend factors did not actually match the non-premultiplied output the
shader itself produced, so a translucent sprite --- one with an alpha value strictly between
fully transparent and fully opaque --- rendered fully opaque instead. This was completely
invisible to the earlier manual screenshot review covering rotation, flips, and filtering; it
was only caught once a real pixel-value assertion replaced ``does this look right in a
screenshot.'' The fix matched Vulkan's own blend-factor pairing. This is the same lesson
Chapter~\ref{ch:spritebatch} drew from SDL\_Renderer's ignored-transform-matrix bug: a
rendering pipeline can look entirely correct in every way a human screenshot review checks,
and still be measurably wrong.

\section{3D depth testing and the gap it exposed}

The first 3D draw path (\texttt{DrawColoredPrimitives} / \texttt{DrawIndexedColoredPrimitives})
was verified with a genuine near/far depth-ordering test rather than trusting draw order
alone. Building it surfaced a real, previously-invisible gap: applying a
\cnaclass{DepthStencilState} had no effect on this renderer at all before this point --- the
apply-state entry point was entirely unimplemented, meaning
\texttt{GraphicsDevice.DepthStencilState} had silently done nothing on WebGPU up to that
point. From there, \cnaclass{PbrEffect} (a real glTF~2.0 metallic-roughness BRDF, ported line
for line from EasyGL's own GLSL implementation) and skinned variants of both
\cnaclass{SkinnedEffect} and \texttt{PbrEffect} (bone-palette skinning up to 72 bones, again
ported directly from EasyGL) were added and both pass their dedicated test suites in full.

\subsection{PbrEffect: a CNAEXT metallic-roughness material}

\cnaclass{PbrEffect} is worth a concrete look before moving on, since --- unlike every stock
effect in Chapter~\ref{ch:stock-effects} --- it has no real-XNA counterpart at all: real XNA
4.0 predates the glTF~2.0 metallic-roughness PBR model by years, so this is a \texttt{CNAEXT}
addition ported line for line from EasyGL's own GLSL implementation, not a port of anything
Microsoft ever shipped. Its property surface mixes ordinary \cnaclass{IEffectMatrices}/%
\cnaclass{IEffectLights} / \cnaclass{IEffectFog} plumbing (identical in shape to every stock
effect) with a genuinely new four-texture-map PBR material:

\begin{lstlisting}
PbrEffect effect(getGraphicsDeviceProperty());
effect.setWorldProperty(worldMatrix);
effect.setViewProperty(camera.View());
effect.setProjectionProperty(camera.Projection());

effect.setTextureProperty(&baseColorMap);               // glTF "baseColorTexture"
effect.setNormalMapProperty(&normalMap);
effect.setMetallicRoughnessMapProperty(&metalRoughMap); // one texture: G=rough, B=metal
effect.setOcclusionMapProperty(&aoMap);
effect.setEmissiveMapProperty(&emissiveMap);

// Both factors multiply the map's channel value (glTF's own "value * factor" convention),
// or stand alone with no map bound:
effect.setMetallicFactorProperty(1.0f);
effect.setRoughnessFactorProperty(1.0f);
effect.setEmissiveFactorProperty(Vector3(0.0f, 0.0f, 0.0f));

effect.setLightingEnabledProperty(true);
effect.EnableDefaultLighting();
\end{lstlisting}

The \texttt{MetallicRoughnessMap} single-texture, two-channel packing above is not an
implementation shortcut CNA introduced --- it is glTF~2.0's own real specified layout for this
material model (green channel roughness, blue channel metallic, red and alpha unused), and
\texttt{PbrEffect} follows it exactly so a metallic-roughness texture exported by any
standards-conformant glTF pipeline can be bound directly with no channel-repacking step. This
is also, concretely, the effect Chapter~\ref{ch:models-meshes}'s two independent skinning
systems both connect to on this renderer: the WebGPU port adds a skinned variant of
\texttt{PbrEffect} alongside skinned \cnaclass{SkinnedEffect}, so a glTF character rig
authored with a metallic-roughness material and driven by \cnaclass{AnimationPlayer} reaches
GPU skinning through this effect rather than through \cnaclass{SkinnedEffect} at all.

\section{Render targets, environment mapping, instancing, and mip generation}

\cnaclass{RenderTarget2D} support combines depth and stencil into a single format regardless
of what was requested (unlike Vulkan's newer per-render-target depth-format selection,
Chapter~\ref{ch:vulkan-backend}),
and needed a real type-safety fix when sampling a render target back as an ordinary texture.
MSAA support, once implemented, initially appeared to fail its own dedicated test --- but
investigation found the infrastructure was already correct, and the reported failure was a
defect in the \emph{test itself} (a missing \texttt{RasterizerState::CullNone}), not in the
MSAA implementation. \cnaclass{EnvironmentMapEffect} and genuine hardware instancing were
added together, with instancing implemented as a real per-instance vertex stream needing no
bind-group changes, unlike the effect families that do require one per draw.
\cnaclass{RenderTargetCube} support surfaced a cross-cutting sprite-geometry bug that has since
been fixed: the original queue path computed clip-space positions from the backbuffer's logical
dimensions even while a differently-sized off-screen target was bound. \texttt{REMED-GFX-019}
now derives dimensions from the current \cnaclass{RenderTarget2D} or cube face instead. Mip
generation for \cnaclass{Texture2D}
and \cnaclass{TextureCube} is implemented as a genuine filtered downsample (a series of
full-screen-triangle render passes, since \texttt{wgpu-native} v29 has no direct
blit-image-equivalent call). Notably, this renderer auto-regenerates the chain after every
level-zero write to a plain \cnaclass{Texture2D} or \cnaclass{TextureCube}. Real XNA/FNA and
the other audited plain-texture paths keep upper levels explicitly authored; SDL\_GPU is the
narrow exception for plain cubes, where a complete level-zero face upload also invokes its
native generator, while its plain 2D and 3D textures remain authored-level resources.
WebGPU's behavior prevents undefined upper levels after a level-zero upload, but it is a
deliberate timing divergence rather than an unqualified superiority: a later level-zero write
can overwrite upper levels that a game authored earlier. Render-target mip generation remains
the separate unbind/pass-finalization contract.

\begin{sourcenote}
The default surface is equally independent of public presentation-format fields, but its
choice is unusually visible. \texttt{ConfigureSurface()} prefers BGRA8-sRGB, then
RGBA8-sRGB, before either UNORM variant. Thus a public \texttt{Color} request can accompany
a native sRGB swapchain, the opposite default-color policy from Vulkan. Its fixed depth
texture uses \texttt{Depth24PlusStencil8}. The renderer recreates it for every non-minimized
surface, including when the public depth value is \texttt{None}. SDL performs any fullscreen window transition;
the next frame detects its physical size and reconfigures the surface, but neither requested
format enum participates.
\end{sourcenote}

\subsection{The corrected render-target-size trap}

The former sprite-geometry bug is worth walking through because it produced no error of any
kind --- only a subtly wrong result:

\begin{lstlisting}
// e.g. a 256x256 render target, while the game window itself is 1920x1080:
device.SetRenderTarget(&smallOffscreenTarget);
spriteBatch.Begin();
spriteBatch.Draw(sourceTexture, Vector2::Zero, Color::White); // "fill the target"
spriteBatch.End();
\end{lstlisting}

Before \texttt{REMED-GFX-019}, the same \texttt{Draw()} call intended for the
$256\times256$ target was converted as though the target were the $1920\times1080$ window,
silently changing its scale and placement. The current \texttt{QueueSprite()} instead selects
the bound 2D target's width/height, the bound cube face's size, or the backbuffer logical
viewport as three explicit cases before converting pixels to NDC.

The dedicated \texttt{WebGPU\_SpriteBatch\_RenderTarget} regression makes the distinction
discriminating rather than visual-only: a $96\times72$ backbuffer is paired with asymmetric
$48\times32$ and $64\times40$ targets; the same destination rectangle must occupy the same
target-relative pixels in both. It also covers target--backbuffer--target isolation, a
non-identity SpriteBatch transform, texture orientation, and a $48\times48$ cube face. Probe
pixels inside the old half-scale rectangle but outside the correct one must remain black. The
old practical prohibition on differently-sized targets is therefore no longer needed.

\subsection{RenderTargetCube and its former silent cast failure}

\cnaclass{RenderTargetCube} support (Task WEBGPU-114) landed after the plain
\cnaclass{RenderTarget2D} path above, and reused most of its design directly: one shared
6-array-layer \texttt{WGPUTexture} (the same layout \texttt{WebGPUTextureCubeRenderer} uses
for an ordinary sampled, readback-capable \cnaclass{TextureCube}), one shared
\texttt{size}$\times$\texttt{size}
depth/stencil texture reused across all six faces (only one face is ever bound at a time, the
same choice Vulkan's own \cnaclass{RenderTargetCube} renderer makes), and a
\texttt{WGPUTextureViewDimension\_2D} view per face for the render-pass color attachment
alongside one \texttt{WGPUTextureViewDimension\_Cube} view spanning all six layers for sampling
the whole thing back later.

That whole-cube sampling view is where a real, previously-invisible bug lived. Before this
task, \cnaclass{EnvironmentMapEffect}'s \texttt{envMap} field on this renderer was hardcoded to
the concrete type \texttt{const WebGPUTextureCubeRenderer*} --- so binding an ordinary
\cnaclass{TextureCube} as \texttt{EnvironmentMapEffect.EnvironmentMap} worked, but
binding a \cnaclass{RenderTargetCube} (a different concrete class implementing the same
\cnaclass{ITextureCubeRenderer} interface) did not: the hardcoded-type cast failed silently,
and the effect fell back to its own 1$\times$1 white cube default instead of throwing or
producing an obviously wrong image. The fix is a small shared interface, read directly from
the real header:

\begin{lstlisting}
// WEBGPU-114: the cube-map sibling of IWebGPUSamplable -- lets both
// WebGPUTextureCubeRenderer and WebGPURenderTargetCubeRenderer
// (render-into) resolve safely to their own whole-cube view wherever a
// GpuDrawParams::envMap needs to be sampled, via a dynamic_cast that
// degrades to nullptr for a null or incompatible-type input.
class IWebGPUCubeSamplable
{
public:
    virtual ~IWebGPUCubeSamplable() = default;
    [[nodiscard]] virtual WGPUTextureView CubeView() const = 0;
};
\end{lstlisting}

Both concrete classes now implement \texttt{IWebGPUCubeSamplable} alongside their own primary
interface --- \texttt{WebGPUTextureCubeRenderer} and \texttt{WebGPURenderTargetCubeRenderer}
alike --- and the draw-dispatch code resolves through it with a single \texttt{dynamic\_cast}
rather than the old hardcoded
concrete-type cast:

\begin{lstlisting}
[[nodiscard]] const IWebGPUCubeSamplable* ResolveCubeSamplable(const ITextureCubeRenderer* tex)
{
    return tex != nullptr ? dynamic_cast<const IWebGPUCubeSamplable*>(tex) : nullptr;
}
// ...
command.envMap = ResolveCubeSamplable(params.envMap);
\end{lstlisting}

\texttt{modules/renderers/webgpu/examples/\allowbreak{}webgpu\_rendertargetcube\_test.cpp}'s own Check C is built specifically to
prove this fix rather than merely exercise the feature: it clears all six faces of a real
\cnaclass{RenderTargetCube} to blue, binds it as \texttt{EnvironmentMapEffect.EnvironmentMap},
draws a full-screen reflective quad, and asserts the center pixel is genuinely blue --- a check
that could not have passed under the old hardcoded-type cast no matter how correctly the rest
of the render-target machinery worked, since the cast itself was the failure point.

The same test file's Check D is worth walking through on its own, because it targets a
different, easy-to-get-wrong question: does switching \emph{into} a cube face and back
correctly restore the backbuffer's own render pass, with nothing leaking either direction?

\begin{lstlisting}
device.SetRenderTarget(nullptr);              // targeting the backbuffer
device.Clear(Color::Blue);
device.SetRenderTarget(&cubeTarget, CubeMapFace::PositiveZ);
device.Clear(Color::Red);                     // targets the cube face, not the backbuffer
device.SetRenderTarget(nullptr);              // back to the backbuffer

// The backbuffer's own centre pixel must STILL be Blue -- the intervening
// cube-face Clear(Red) must not have leaked into it -- and the cube face
// itself must genuinely have received Red, not been silently skipped.
\end{lstlisting}

This mirrors the identical architecture check Task WEBGPU-53/54 already established for plain
\cnaclass{RenderTarget2D} in the section above, extended to confirm the same "eager flush on
target switch" design generalizes correctly to a third target kind rather than being
coincidentally correct for the two it was originally built against. Two scope cuts are
documented honestly rather than silently: mip-chain regeneration throws for
\cnaclass{RenderTargetCube} (\texttt{mipMap=true} is rejected, matching plain
\cnaclass{RenderTarget2D}'s own precedent on this renderer), and MSAA is not implemented ---
\texttt{MultiSampleCount} is accepted and ignored, always reporting \texttt{0} honestly rather
than silently pretending to honor a request it cannot fulfil.

\subsection{Why block-compressed texture upload remains open}

The ``no renderer anywhere in the project does real block-compressed texture upload'' gap named
below is worth a fuller account on this renderer specifically, because it was genuinely
investigated here rather than assumed unfixable. A direct diagnostic against this project's own
pinned \texttt{wgpu-native} v29.0.1.1, run on the real development GPU, confirmed the hardware
side is not the obstacle at all: \texttt{wgpuAdapterHasFeature(adapter,
TextureCompressionBC)} returned true, and requesting a device with that feature in
\texttt{requiredFeatures} succeeded too --- this machine's real adapter genuinely supports
DXT1/3/5 upload today. The actual blocker sits one layer up, and applies project-wide, not just
to \texttt{WEBGPU}: \texttt{Texture2D::FromStream()} always fully CPU-decompresses DXT-compressed
source bytes to plain RGBA8 \emph{before} any renderer ever sees them, and the shared
\texttt{CNA::Internal::Graphics::ImageData} struct every renderer's texture-creation path
consumes has no field at all for a surface format or a compressed-bytes flag --- it is
hardcoded to a plain \texttt{std::vector<uint8\_t>} of RGBA8 pixels. A WebGPU-only workaround
was considered and deliberately rejected: hand-crafting a direct \texttt{IGraphicsRenderer::%
CreateTexture()} call with raw BC1 bytes stuffed into that RGBA8-typed field would compile and
might even render correctly on this one renderer, but would be reachable by no real game (nothing
in the actual \cnaclass{Texture2D} loading path could ever produce it) and would misrepresent
\texttt{ImageData}'s own documented contract --- exactly the shape of shortcut this project's
own culture treats as worse than leaving the gap honestly open. Closing this for real needs a
cross-renderer \texttt{ImageData} design change (or a parallel compressed-texture entry point)
plus a matching \texttt{Texture2D.cpp} change to stop always-decompressing DXT source data ---
deliberately scoped out of any single renderer's own task list, tracked as a dedicated follow-up
instead.

\section{Three former limitations that are no longer true}

The WebGPU plan moved quickly after this chapter's first draft, so three earlier shorthand
limitations need a precise correction rather than being allowed to turn into folklore. First,
an arbitrary plain \cnaclass{Texture2D} now has a real renderer-level \texttt{GetData()} path.
That method on \texttt{WebGPUTextureRenderer} copies the requested mip level into a temporary
\texttt{MAP\_READ} buffer with WebGPU's required row alignment, waits through the asynchronous
map callback, and then extracts the requested $x,y,w,h$ rectangle. The dedicated
\texttt{WebGPU\_Texture2D\_GetData} test proves three independently useful cases: an exact
full-gradient round trip, a non-origin $2\times2$ sub-rectangle, and a distinct level-one mip
round trip. That is deliberately tested through \texttt{IGraphicsRenderer}, not merely through
the public \cnaclass{Texture2D}: the latter normally returns its shared CPU pixel shadow for a
plain texture, so an XNA-layer round trip alone could pass while this GPU copy path remained a
silent base-class no-op.

Second, culling, scissor, and viewport state are no longer merely stored values. Culling is
baked into the WebGPU 3D pipeline key and mapped to \texttt{WGPUCullMode}; a differential
pixel test proves that a deliberately chosen winding disappears under one XNA cull mode and
remains visible under the other. Scissor and viewport are genuine dynamic render-pass state via
\texttt{wgpuRenderPassEncoderSetScissorRect} and
\texttt{wgpuRenderPassEncoderSetViewport}. Wiring the latter exposed an otherwise-hidden
logical-versus-physical-size hazard: CNA's stored logical viewport can exceed the real
swapchain dimensions. The renderer therefore clamps it at application time, which is both a
WebGPU API requirement and what restored the pre-existing 2D and 3D regression tests.
Every queued SpriteBatch or 3D command captures its own viewport and scissor state at the public
draw call. Replay applies those snapshots dynamically through the two native setters on the
backbuffer, 2D-target, and cube-face paths; consecutive commands with identical state skip the
redundant native call. Thus several changes inside one deferred pass no longer collapse to the
last live value. The shared deferred pixel fixtures prove the per-command result, while the
WebGPU cardinality tests prove it did not require a pipeline variant, render-pass split, extra
submit, or one setter call for every same-state draw. A disabled or zero scissor expands to the
full current target instead of issuing an invalid zero-area rectangle.

Third, this is not a claim that every state object is complete. \cnaclass{FillMode}'s
\texttt{WireFrame} value is stored but polygon draws now refuse it because the pinned WebGPU API
offers no polygon-mode switch. The renderer stores the complete stencil configuration but has
deliberately not yet baked stencil operations
into its many pipeline families. The last deferral is evidence-based rather than cosmetic:
the cull-mode mapping itself required a pixel test to correct an initially plausible but wrong
winding interpretation, so copying a stencil front/back mapping without an equivalent
differential stencil test would risk making an invisible wrongness look ``implemented.''

\section{Clear ordering and usage are consistent across target types}

Each public clear enters the same ordered stream as all eleven deferred draw families, carrying
independent colour, depth, and stencil flags and values. Replay partitions one bind cycle into
native pass segments: leading/consecutive clears can fold into the next segment's load actions,
while a clear after a draw closes that segment so the following pass observes it at the exact
public position. A discard-content target may clear its first segment as its usage permits;
later explicit clears remain ordered rather than becoming one global final value. The shared
clear-options and ordered-clear fixtures cover backbuffer, 2D-target, cube-face, and MRT-rejection
boundaries without an intervening flush.

The shared FNA predicate maps only \texttt{DiscardContents} to discard; Preserve and Platform
both preserve. WebGPU receives that Boolean for 2D and cube targets, and the first native
segment selects clear versus load for colour, depth, and stencil together. Later segments load
prior contents unless an explicit ordered clear selects an aspect. A cube face owns its colour
view while all faces share the target's depth/stencil attachment, matching the reference
contract. Cube-usage and depth/stencil-usage regressions distinguish all three enums and the
six faces rather than inferring behavior from a combined clear.

The backbuffer's
\texttt{Presentation\allowbreak{}Parameters.Render\allowbreak{}Target\allowbreak{}Usage} remains
outside this off-screen policy, and plural binding still rejects MRT rather than creating more
than one attachment. The full comparison is \S\ref{sec:backend-clear-matrix}; the separate
usage matrix is \S\ref{sec:backend-rendertarget-usage-matrix}.

\section{What remains open}

Several gaps are shared with every other CNA renderer rather than unique to this one: no
renderer anywhere in the project does real block-compressed texture upload (the paragraph above
covers exactly why, and why it is a shared architectural gap rather than a per-renderer one),
and the cross-renderer design change that would make it reachable from a real game remains
larger than any one renderer task. Specific to this renderer: multiple simultaneous render
targets, actual stencil-operation mapping (not merely the already-stored state and dynamic
reference value), \texttt{WireFrame} rendering,
custom \cnaclass{SpriteBatch} effects, and custom WGSL effects authored outside the stock-effect
set. \cnaclass{BlendState} is no longer on that list for the one attachment WebGPU can bind:
all source/destination factors and operations, dynamic \texttt{BlendFactor}, slot-zero
\texttt{ColorWriteChannels}, and \texttt{MultiSampleMask} reach keyed SpriteBatch and 3D
pipelines with dedicated differential tests. Slots 1--3 remain inapplicable until MRT exists.
The two \cnaclass{RenderTargetCube} scope cuts also remain honest: mipmapped cube targets
throw at construction, and their per-target MSAA request is accepted but reports zero rather
than claiming an unimplemented multisampled array-texture resolve. Browser/Emscripten WebGPU is
entirely out of scope so far --- this renderer is native \texttt{wgpu-native} only.

\section{Verification methodology, and how it differs from D3D9's}

This renderer's own verification runs on a native automated CTest suite (gracefully skipping
where no Wayland/X11 display is available) plus a battery of dedicated pixel-assertion tests
per feature area, and the one independent real-game integration test described above. It is
explicitly \emph{not} covered by the oracle-versus-real-XNA methodology
Chapter~\ref{ch:direct3d-backends} describes for \texttt{D3D9} / \texttt{D3D11} / \texttt{D3D12}
--- that comparison against a genuine, running XNA~4.0 reference is reserved for the three
Direct3D renderers specifically, not applied here.
